\documentclass[12pt]{article}
\usepackage[utf8]{inputenc}
\usepackage[T1]{fontenc}
\usepackage{hyperref}
\title{Magori Language Resources for K-12 Curriculum Development}
\author{Compiled by Claude Sonnet 5.5 (claude-sonnet-5-5) for Andr\'as Kornai}
\date{}
\begin{document}
\maketitle

\section{Introduction}
Magori (ISO 639-3 zgr; Glottocode mago1248) is a small Austronesian language of Central Province, Papua New Guinea. It is spoken at the eastern end of Table Bay, on the lower Bailebo-Tavenei River. Wikipedia places it in the Oumic branch of Papuan Tip Oceanic and lists about 100 speakers (2000); UNESCO classifies it as Vulnerable. Many speakers are bilingual or trilingual in Suau, Mailu or Hiri Motu. Its vocabulary shows heavy influence from Mailu. The searches found only reference and wordlist material, and no educational materials.

\section{Primary Educational Resources}
No primers, textbooks or school programmes in Magori were found.

\section{Linguistic and Reference Materials}
\begin{itemize}
\item Wikipedia, ``Magori language''. \url{https://en.wikipedia.org/wiki/Magori_language}
\item Endangered Languages Project entry for Magori. \url{https://endangeredlanguages.com/lang/2267}
\item Verbix page on Magori. \url{https://docs.verbix.com/Languages/Magori}
\item ASJP (Automated Similarity Judgment Program) wordlist entry for Magori. \url{https://asjp.clld.org/languages/MAGORI}
\end{itemize}

\section{Religious and Community Resources}
None were found in the search results.

\section{Digital and Audio Resources}
No audio or video resources were found in the search results. The Endangered Languages Project entry above may list further sources.

\section{Curriculum Development Context}
Almost nothing exists for educational use. A curriculum would begin with community-led documentation of vocabulary, texts and recordings, and would take account of the multilingual setting (Suau, Mailu and Hiri Motu).

\end{document}